
\begin{table*}[ht]
\centering
\begin{minipage}{0.45\textwidth}
\centering
\setlength{\tabcolsep}{1mm}
\begin{tabular}{lrrr}
\toprule
\textbf{ANWR TX Cranes} & \textbf{MAE} & \textbf{RMSE} & \textbf{BPR-50} \\
\midrule
Last timestep & 0.24 & 1.04 & 0.27 \\
Avg over 10 & 0.25 & 0.78 & 0.38 \\
Chance decision, $\hat{y}=0$ & 0.15 & 0.79 & 0.05 \\
EpiGNN & 0.38 & 0.70 & 0.06 \\
PG & 1.03 & 4.85 & 0.35 \\
SPO+ & 0.15 & 0.83 & 0.18 \\
NLL Only & 0.35 & 1.14 & 0.38 \\
DAML ($\epsilon=1$) & 9.08 & 35.89 & 0.39 \\
BPR Only & 40495 & 40495 & 0.41 \\
\bottomrule
\end{tabular}
\caption*{(a) ANWR TX Cranes 2009-2010.}
\end{minipage}
%\hfill
~~~~
\begin{minipage}{0.45\textwidth}
\centering
\setlength{\tabcolsep}{1mm}
\begin{tabular}{lrrr}
\toprule
\textbf{Cook County IL Overdose} & \textbf{MAE} & \textbf{RMSE} & \textbf{BPR-100} \\
\midrule
Last timestep & 1.07 & 1.66 & 0.76 \\
Avg over 5 & 0.99 & 1.58 & 0.80 \\
Chance decision, $\hat{y}=0$ & 1.37 & 2.46 & 0.20 \\
EpiGNN & 1.33 & 2.03 & 0.32 \\
PG & 5.65 & 17.03 & 0.80 \\
SPO+ & 1.25 & 2.38 & 0.74 \\
NLL Only & 1.03 & 1.58 & 0.78 \\
DAML ($\epsilon=1$) & 1.12 & 1.84 & 0.80 \\
BPR Only & 70.31 & 152.72 & 0.82 \\
\bottomrule
\end{tabular}
\caption*{(b) Cook County IL 2021-2022.}
\end{minipage}
\caption{Performance comparison between different model families for error-based and decision metrics. The top 2 rows represent simple historical baselines: using either the previous timestep alone, or an average over a larger amount (10 bi-months for the crane dataset, and 5 years for Cook County). Next, the \emph{Chance decision} model chooses spatial locations by random chance, and uses all 0's for predicted $\hat{y}$ in MAE and RMSE calculations. Despite the simplicity, this is a competitive model on the error-based metrics due to sparsity. It outperforms all others on MAE on the Cranes dataset. Next is EpiGNN, a spatiotemporal forecasting model based on graph neural networks, trained to minimize RMSE. Although this model has the best performing RMSE on the crane dataset, it makes poor top-K decisions on both. Finally, the last three rows show different ways to train the negative binomial mixed effects regression model in \ref{eq:model}. \emph{NLL Only} trains to maximize likelihood alone (\ref{eq:ml_obj}), \emph{BPR Only} is our direct-loss minimization trained to maximize BPR alone (\ref{eq:dlm_obj}), and DAML (\ref{eq:hybrid_obj}) is our hybrid loss that allows a user to explore the pareto frontier between likelihood and BPR-K.}
\label{tab:results}
\end{table*}
